% misc/custom_defs.tex
% Defining all custom commands in a dedicated file helps you to keep the main
% file clean. You can also load and configure additional packages there.

%-------------------------------------
%       Additional packages
%-------------------------------------
% Here you can load other packages.

\usepackage{xspace}             % smart space after custom commands
% \usepackage[option1,option2]{coolpackage}


%-------------------------------------
%       Additional commands
%-------------------------------------
% Here you can define your own commands.

%%%%% Logical markup
% Commands to style logical elements of your text. Adapt to your preferences!
\renewcommand{\titlegraphic}[1]{\vbox{\centering #1}}
\newcommand{\species}[1]{\textit{#1}}       % species names in italics
\newcommand{\seq}[1]{\texttt{#1}}           % RNA/DNA sequence
\newcommand{\file}[1]{\texttt{#1}}          % file names
\renewcommand{\cmd}[1]{\texttt{#1}}         % executable code and commands
\newcommand{\prog}[1]{\textit{#1}}          % names of programs & packages
\newcommand{\email}[1]{\href{mailto:#1}{\texttt{#1}}} % email with mailto link
\newcommand{\quo}[1]{\enquote{#1}}          % quotation marks, dep. on language
\newcommand{\tblhead}[1]{\textbf{#1}}       % header of tables
\newcommand{\pkg}[1]{\texttt{#1}}           % LaTeX package
\newcommand{\clsopt}[1]{\texttt{#1}}        % LaTeX class option


%%%%% Abbreviations
% Common abbreviations that are never spelled out. For domain-specific
% abbreviations that should be added to the List of Abbreviations, use the
% \defabb command instead!
\newcommand{\vs}{\emph{vs.\@}\xspace}   % vs as in versus
\newcommand{\eg}{e.\,g.\@\xspace}       % e.g. (for example)
\newcommand{\zb}{z.\,B.\@\xspace}       % z.B. (zum Beispiel)
\newcommand{\ie}{i.\,e.\@\xspace}       % i.e. (that is)
\newcommand{\wrt}{w.\,r.\,t.\@\xspace}  % w.r.t. (with respect to)
\newcommand{\cf}{cf.\@\xspace}          % cf. (see/refer to)
\newcommand{\fivepr}{\ensuremath{5^\prime}\xspace}      % 5' (five prime)
\newcommand{\threepr}{\ensuremath{3^\prime}\xspace}     % 3' (three prime)
\defabb{ToC}{Table of Contents}
\defabb{UTC}{Coordinated Universal Time}
\defabb{DAG}{Directed Acyclic Graph}
\defabb{mFAS}{minimum Feedback Arc Set}
\defabb{SV}{Super Vertex}[Super Vertices]

%%%%% Math
% For custom operators and math commands.
% \DeclareMathOperator{\var}{Var}             % variance
% \DeclareMathOperator*{\argmin}{arg\,min}    % arg min


%-------------------------------------
%       Symbols & Categories
%-------------------------------------
% Here you can place your own symbols and categories, cf. the manual.

%%%%% Categories
\addSymbolCategory{C}{Cycle}
\addSymbolCategory{T}{Tree}

%%%%% Symbols
% \sym{C}{The circumference of a cycle.}
% \sym{\pi}[kreis]{The ratio of a circle's circumference to its diameter.}[C]

% End of misc/custom_defs.tex
